\section{How much recent history is enough?}
\label{sec:capability}
\subsection{Capability recovery and its price}
State records selected facts; recent history also retains the wording,
ordering, and surrounding context of interactions. Adding it raises
model-written success from \ModelStateRate\% at $k=0$ to
\ModelHistoryTwoRate\% at $k=2$ and \ModelHistorySixRate\% at $k=6$.
The last setting differs from full history by
$\ModelHistorySixDelta$ percentage points, with interval
$\ModelHistorySixDeltaCI$. This is substantial recovery, not demonstrated
equivalence.

\begin{table}[htbp]
\centering\small
\caption{Recent-history retention, quality and economics. All retained
outcomes contribute. Each memory family has its own full-history control.
Cost ratio is the median matched ratio, not a ratio of medians; unavailable
hybrid ratios are left blank. Solved/\$ uses aggregate cost, including failures.}
\label{tab:history}
\resizebox{\linewidth}{!}{\input{tables/new-history}}
\end{table}

\begin{table}[htbp]
\centering\small
\caption{Resource panel for the same history policies. Peak prompt and
token counts are execution-level medians; cache hit is the median
within-execution cached fraction. Hybrid peak and cache-hit statistics
are unavailable, not zero.}
\label{tab:history-resources}
\resizebox{\linewidth}{!}{\input{tables/new-history-resources}}
\end{table}

The hybrid policy shows the same direction: \HybridRate\% without history,
\HybridHistoryTwoRate\% at $k=2$, and \HybridHistorySixRate\% at $k=6$.
The latter improves on hybrid state alone by
$\HybridHistorySixVsHybridDelta$ percentage points
($p=\HybridHistorySixVsHybridP$), against a full-history reference of
\HybridControlRate\%.
Its cost per observed success increases from \$\HybridCostPerSuccess{}
to \$\HybridHistorySixCostPerSuccess{}.
For model-written memory it instead falls from
\$\ModelStateCostPerSuccess{} to \$\ModelHistorySixCostPerSuccess{}.
The same retention decision therefore has different economic consequences.

% Adapted from the repository's original trade-off scatter.
\begin{figure}[tbp]
\centering
\begin{tikzpicture}
\begin{axis}[
width=.47\linewidth,height=5.1cm,
title={Model-written memory},title style={font=\small},
xlabel={Mean cost (\$)},ylabel={Solved rate (\%)},
xmin=.040,xmax=.078,ymin=30,ymax=68,
xtick={.04,.05,.06,.07},scaled x ticks=false,
xticklabel style={/pgf/number format/fixed,/pgf/number format/precision=2,font=\scriptsize},
yticklabel style={font=\scriptsize},label style={font=\small},
grid=major,grid style={black!12},
nodes near coords={\pgfplotspointmeta},point meta=explicit symbolic,
every node near coord/.append style={font=\scriptsize,anchor=south,yshift=3pt}]
\addplot[only marks,mark=*,mark size=2.5pt,blue!65!black]
table[col sep=comma,x=mean_cost_usd,y=solved_pct,meta=history]{../paper/generated/new-model-history.csv};
\end{axis}
\begin{axis}[
at={(.52\linewidth,0)},anchor=south west,
width=.47\linewidth,height=5.1cm,
title={Hybrid memory},title style={font=\small},
xlabel={Mean cost (\$)},
xmin=.016,xmax=.028,ymin=30,ymax=68,
xtick={.016,.020,.024,.028},scaled x ticks=false,
xticklabel style={/pgf/number format/fixed,/pgf/number format/precision=3,font=\scriptsize},
yticklabel style={font=\scriptsize},label style={font=\small},
grid=major,grid style={black!12}]
\addplot[only marks,mark=square*,mark size=2.5pt,green!45!black]
table[col sep=comma,x=mean_cost_usd,y=solved_pct,meta=history]{../paper/generated/new-hybrid-history.csv};
\node[font=\scriptsize,green!45!black,anchor=south west,xshift=3pt,yshift=3pt]
 at (axis cs:\HybridControlCostMean,\HybridControlRate) {full};
\node[font=\scriptsize,green!45!black,anchor=north east,xshift=-3pt,yshift=-3pt]
 at (axis cs:\HybridHistorySixCostMean,\HybridHistorySixRate) {6};
\node[font=\scriptsize,green!45!black,anchor=north east,xshift=-3pt,yshift=-3pt]
 at (axis cs:\HybridCostMean,\HybridRate) {0};
\node[font=\scriptsize,green!45!black,anchor=south west,xshift=3pt,yshift=3pt]
 at (axis cs:\HybridHistoryTwoCostMean,\HybridHistoryTwoRate) {2};
\end{axis}
\end{tikzpicture}
\caption{Success--cost trade-offs when retaining history. Labels are the retention setting $k$; full denotes the full-history reference without a state block. Each panel includes failures in mean cost and uses its own control and cost convention. The horizontal scales differ.}
\label{fig:pareto}
\end{figure}


\subsection{The complete window sweep and why six}
Recovery at $k=6$ alone cannot show that a larger window adds nothing.
The pilot therefore also ran $k=12$ and $k=\infty$ (state plus the full
transcript, distinct from the no-state control), with three repetitions
on \WindowTaskCount{} tasks (Tables~\ref{tab:window}
and~\ref{tab:window-tasks}).

\begin{table}[htbp]
\centering\small
\caption{Window pilot in the standard two panels (finalized medians over
three repetitions). Per-execution solved rates, efficiencies, pooled cost
ratios, and input counters were not retained (--). Cache hit is the
reported pilot value, not reconstructed from counters.}
\label{tab:window}
\resizebox{\linewidth}{!}{\input{tables/new-window-quality}}

\medskip
\resizebox{\linewidth}{!}{\input{tables/new-window-resources}}
\end{table}

\begin{table}[htbp]
\centering\small
\caption{Window pilot outcomes reported in the research log (raw sessions
unavailable). A task counts as solved if any of its three repetitions
passed. Discordant tasks and exact McNemar $p$ compare with full history.
Output per turn is the reported median.}
\label{tab:window-tasks}
\resizebox{\linewidth}{!}{\input{tables/new-window-tasks}}
\end{table}

Three measurements stop changing at $k=6$. Tasks solved rise from
\WindowStateTasksSolved{} of \WindowTaskCount{} at $k\leq2$ to
\WindowSixTasksSolved{} at $k=6$ and stay there; the repetition ratio
falls from \WindowStateRepetition{} to \WindowSixRepetition{} and stays
there; and the $k=6$, $12$, and $\infty$ rows each differ from the control
by a single discordant task. Two of these flat lines are bounds, not
plateaus. Solving all \WindowTaskCount{} tasks is the ceiling of the
task set, and a ratio of one is the floor of the repetition diagnostic. A 12-task pilot
needs at least six one-sided discordant tasks for significance at the five percent level, so it cannot
separate conditions that are all at or near the ceiling. The honest
reading is that the threshold lies between two and six messages, not that
wider windows are useless.

The resource columns keep moving. Median peak prompt grows from
\WindowSixPeakMedian{} thousand tokens at $k=6$ to \WindowTwelvePeakMedian{}
and \WindowInfinitePeakMedian{}; median turns fall from
\WindowSixCallsMedian{} to \WindowTwelveCallsMedian{} and
\WindowInfiniteCallsMedian{}; median cost is \$\WindowSixCostMedian{},
\$\WindowTwelveCostMedian{}, and \$\WindowInfiniteCostMedian{}. Matched
cost ratios stay between \WindowInfiniteCostRatioMin{} and
\WindowTwelveCostRatioMax{} for every $k\geq6$ in every repetition
(Supplementary~\ref{app:window}). The reason is visible in output per turn.
Every state condition writes about twice the control's
\WindowControlOutputPerTurn{} tokens (\WindowSixOutputPerTurn{} at $k=6$),
because the model rewrites state on each step. Widening the window cannot
remove that fixed overhead; it only trades fewer turns for larger prompts.
Cache hit also stays low (\WindowSixCacheHit--\WindowInfiniteCacheHit\%
against \WindowControlCacheHit\%), consistent with a state block that
precedes the transcript in Eq.~\eqref{eq:ctx-state} and changes every turn.

So $k=6$ is the smallest tested window at which success and redundant
re-invocation reach the pilot's resolution limit, and the cheapest point
beyond it is not identified by these data. The main model-state comparison
confirms recovery at $k=6$ but did not run $k=12$ or $k=\infty$; a
discriminating follow-up needs harder tasks with headroom above the
ceiling and a no-state recent-history control.
